% ===========================================================
\section{Introduction}
% ===========================================================

MAX phases are a family of ternary layered ceramics with the general
formula M\textsubscript{2}AX, where M is an early transition metal,
A is an A-group element, and X is either carbon or nitrogen~\cite{Barsoum2000}.
Their unique combination of metallic electrical and thermal conductivity with
ceramic hardness and high-temperature stability makes them attractive
for applications in nuclear cladding, high-temperature coatings,
and next-generation electronics~\cite{Gogotsi2019}.

The lattice-dynamical properties of MAX phases --- encoded in their
phonon dispersion relations $\omega(\vq)$ --- are central to understanding
thermal transport, mechanical stability, and superconductivity.
Conventional \textit{ab initio} calculations based on density functional
perturbation theory (DFPT) are highly accurate but scale poorly:
a single phonon band structure requires tens of CPU-hours,
making high-throughput exploration of compositional space
computationally prohibitive.

Machine-learning interatomic potentials (MLIPs) have emerged as a
powerful tool to bridge this gap~\cite{Behler2007,Batzner2022}.
Graph neural networks (GNNs), in particular, are well suited to
crystalline systems because they operate directly on the atomic graph
and can naturally encode permutation and translational invariance~\cite{Xie2018}.
Recent equivariant architectures such as NequIP~\cite{Batzner2022},
MACE~\cite{Batatia2022}, and eSEN~\cite{TODO} additionally respect
rotational symmetry (E(3)-equivariance),
leading to improved data efficiency and physical consistency.

In this work we make the following contributions:
\begin{enumerate}
  \item We curate a dataset of \textbf{358 MAX-phase compounds} with
        DFT-computed force constants, phonon band structures, and
        elastic constants.
  \item We develop an \textbf{E(3)-equivariant GNN} with radial basis
        function (RBF) edge encoding that predicts full phonon
        dispersion curves from crystal structure alone.
  \item We incorporate \textbf{physics-informed constraints} ---
        IFC symmetry, the acoustic sum rule, and band smoothness ---
        directly into the training objective.
  \item We perform a systematic \textbf{feature importance analysis}
        identifying the atomic descriptors most critical for
        phonon prediction accuracy.
\end{enumerate}
